\scanpage{656}\mspage{44}
On sait bien que l'existence de ces foncteurs,
pour $u$ donné, fait appel à l'existence de
certaines $\varinjlim$ resp. $\varprojlim$ dans $M$ ; p.ex. si
$I'=e$, il est néc. et suffisant que les $\varinjlim$
(resp. $\varprojlim$) de type $I^\circ$ soient \unsure{repr.} dans $M$.
Mais avec les hypothèses faites sur $M$, il
n'y a pas de problème si $u:I\unsure{-|\to}I'$ est
un foncteur entre catégories \uline{finies}.

On va donner un exemple où il n'en
est plus ainsi.

\uline{Exemple 3.11} Soit $k$ un anneau
commutatif, et \deleted{\unsure{$\mathscr D$|$D$}} $M=K^b(k\text{-}\mathrm{Ab})$
la catégorie des complexes de \unsure{cochaines|cochaînes}
bornés, de $k$-modules, $W$ \unsure{les|la} \deleted{\unsure{constituée des}}
quasi-isomorphismes, donc
\[
 \mathrm{Ho}_W=\unsure{\mathrm{Der}|\mathrm{D}\deleted{\mathrm{er}}}^{b}(k)
\]
catégorie dérivée bornée de la catégorie
abélienne des $k$-modules. On sait,
en utilisant l'existence d'assez de
projectifs dans $k\text{-}\mathrm{Ab}$, que cette
structure est associée à une structure
de Quillen \inserted{close} $(W,C,F)$, décrite par Quillen
% archive sheet number: 655
